% Ad Familiares 2.4 --- headnote
% ad-familiares-02-04, 53 BC, Rome

\textit{Cicero to C. Scribonius Curio, written from Rome in 53
BC (the manuscripts give only the year: \textit{Scr. Romae a. 701}).
Curio --- the future tribune of 50 BC whose veto would
become the trigger of the civil war --- is at this date a very
young man, abroad and full of himself. The letter opens the
small sub-correspondence with him preserved in book 2 of the
\textit{Familiares} (2.1--7), the keynote of which is patient,
flattering management of a brilliant patrician's vanity. The
hand is light: a survey of the kinds of letter one might write,
worked into an apology for sending none of them. Of the playful
register Cicero says, by Hercules, no citizen can laugh in these
times; of the serious, what is there that he could write
weightily to Curio except the commonwealth --- and on the
commonwealth he dares not write what he thinks and will not
write what he does not. The shape is the message: a letter
that explains, with elaborate courtesy, why it has nothing in it.}

\textit{The second section is the real one. Cicero falls back
on what he calls his \textit{clausula} --- his standard sign-off
to Curio --- the exhortation to the highest praise. The figure
is forensic: a great expectation has been set up against him
as an adversary, and the way to defeat it is to put in the
labour in the arts that win the distinctions Curio has come to
love. Cicero closes by insisting the prod is not meant to
inflame --- only to testify to his love. The reader knows
already, in 53 BC, that the love is provisional and the
flattery interested; the letter is a piece of careful
political husbandry of a young man who would one day be
useful, or dangerous, or both.}
